\section*{General Introduction}
\addcontentsline{toc}{chapter}{General Introduction}

Since I was young, I've always been fascinated by fluid mechanics.
I remember watching feathers, birds and planes flying and being amazed by the way they could fly.
I was also interested in the way water flows in rivers and sailing boats. In high school, I began working on plane-related projects and explored them more theoretically through \ac{CFD}.

This fascination led me to study physics and mathematics at university, where I discovered applied mathematics including CFD, \ac{AI} and Kalman Filter.
During my studies, I've become passionated about applied mathematics and fluid mechanics. That's why I did my last internship at the \ac{CERFACS}.
It convinced me to do a thesis in CFD after \ac{IPSA}, my current engeneering school.

The subject proposed by the \ac{LMFL} really suited me: CFD using an open-source library, coupled with AI and experimental data. So it was with enthusiasm that I started this internship on February 26, 2025, in this laboratory for a duration of six months.

The main objective of this internship is to improve the numerical simulation of a centrifugal pump using experimental data through \ac{DA}. This allows to improve the understanding of the physical phenomena involved in the pump's operation and then to optimize its performance.

The LMFL is a public laboratory with about 40 researchers and 25 PhD students, spread over three sites located in Lille (\ac{ENSAM} and \ac{ONERA}) and on the Villeneuve d'Ascq campus (Centrale Lille, \ac{CNRS} and University of Lille).

My main tasks during this internship were:

\vspace{-0.15cm}

\begin{itemize}\setlength\itemsep{0.1em}
    \item Explore and understand the limitations of the baseline \ac{IBM};
    \item Establish the set of sensor signals to be employed as high-fidelity inputs in the DA;
    \item Ensure the convergence of the CFD solution using IBM;
    \item Implemente the DA problem in \ac{OpenFOAM} and \ac{CONES};
    \item Optimize the forcing parameters of the IBM using DA with CONES to correct a non-physical behavior.
\end{itemize}

%\vspace{-0.05cm}

To present the details of my internship work, the manuscript is structured as follows:

\autoref{chap1} introduces the LMFL laboratory and the ENSAM Lille campus, with a focus on the scientific environment and the social/environmental 
responsibility of the host institution. 

\autoref{chap2} presents the context of the study: first the state of the art and the problem investigation, then the pre-processing of the experimental data, and finally the numerical ingredients employed, including the CFD governing equations, the OpenFOAM solver, the EnKF method and the CONES library. 

\autoref{chap3} is dedicated to the application to the centrifugal pump case: the modelling of the pump, the preliminary IBM results, the first simple application of the DA, and finally the analysis of the obtained results. 

The report ends with a conclusion, an appendix section gathering additional technical material, the bibliography, and a list of acronyms. 


\begin{comment}
        \item Introduction to DA problems;
        \item Numerical ingredients of the CFD and DA problems;
        \item Preliminary results using IBM;
        \item Application of DA to the IBM;
        \item Results and perspectives.
    \end{enumerate}
\end{comment}
